\documentclass[12pt,a4paper]{article}
\usepackage[utf8]{inputenc}
\usepackage[T1]{fontenc}
\usepackage{amsmath}
\usepackage{listings}
\usepackage{xcolor}
\usepackage{booktabs}
\usepackage[margin=2cm]{geometry}
\usepackage{enumitem}
\usepackage{mdframed}
\usepackage{array}
\usepackage{etoolbox}

% --- Zmniejszenie odstępów ---
\setlength{\parskip}{0.5\baselineskip}
\setlength{\abovecaptionskip}{2pt}
\setlength{\belowcaptionskip}{2pt}

% --- Modyfikacja środowisk mdframed ---
\newmdenv[backgroundcolor=lightyellow, linecolor=orange!60!black, linewidth=1.2pt,
innertopmargin=4pt, innerbottommargin=4pt, innerleftmargin=6pt, innerrightmargin=6pt,
skipabove=4pt, skipbelow=4pt]{cwiczeniebox}

\newmdenv[backgroundcolor=lightgreen, linecolor=green!50!black, linewidth=1.2pt,
innertopmargin=4pt, innerbottommargin=4pt, innerleftmargin=6pt, innerrightmargin=6pt,
skipabove=4pt, skipbelow=4pt]{rozwiazaniebox}

\newmdenv[backgroundcolor=lightblue, linecolor=headerblue, linewidth=1.5pt,
innertopmargin=6pt, innerbottommargin=6pt, innerleftmargin=8pt, innerrightmargin=8pt,
skipabove=4pt, skipbelow=4pt]{infobox}

% --- Ustawienia listings (wieksza czcionka) ---
\lstdefinestyle{matlabstyle}{
	language=Matlab,
	backgroundcolor=\color{backcolour},
	commentstyle=\color{codegreen},
	keywordstyle=\color{codeblue}\bfseries,
	numberstyle=\tiny\color{codegray},
	stringstyle=\color{codegreen},
	basicstyle=\ttfamily\small,
	breakatwhitespace=false,
	breaklines=true,
	captionpos=b,
	keepspaces=true,
	numbers=none,
	showspaces=false,
	showstringspaces=false,
	showtabs=false,
	tabsize=2,
	frame=single,
	rulecolor=\color{codegray},
	aboveskip=4pt,
	belowskip=4pt,
	lineskip=0pt,
	fontadjust=true
}

\lstset{style=matlabstyle}

% --- Zmniejszenie rozmiaru naglowkow sekcji ---
\usepackage{titlesec}
\titleformat{\section}
{\normalfont\Large\bfseries\color{headerblue}}
{}{0em}{}
\titlespacing*{\section}{0pt}{12pt}{4pt}

% --- Kolory ---
\definecolor{codegreen}{rgb}{0,0.55,0}
\definecolor{codegray}{rgb}{0.5,0.5,0.5}
\definecolor{codeblue}{rgb}{0,0,0.85}
\definecolor{backcolour}{rgb}{0.96,0.96,0.93}
\definecolor{headerblue}{rgb}{0.18,0.37,0.60}
\definecolor{lightblue}{rgb}{0.88,0.93,0.98}
\definecolor{lightyellow}{rgb}{1.0,0.98,0.88}
\definecolor{lightgreen}{rgb}{0.90,0.97,0.90}
\definecolor{lightred}{rgb}{1.0,0.93,0.93}

% ---
\title{\textbf{\Large Cwiczenia z MATLAB -- Interpolacja i aproksymacja}\\[4pt]}
\author{}
\date{}

\begin{document}
	
	\maketitle
	\thispagestyle{empty}
	
	
	
	\newpage
	% ----------------------------------------------------------------
	\begin{infobox}
		
		\medskip
		\textbf{Funkcje i instrukcje uzywane na tych zajeciach:}
		\begin{center}
			\renewcommand{\arraystretch}{1.1}
			\begin{tabular}{p{5cm}p{8cm}}
				\toprule
				\textbf{Instrukcja / Funkcja} & \textbf{Do czego sluzy} \\
				\midrule
				\texttt{if / elseif / else / end} & Decyzja: co zrobic gdy warunek jest spelniony \\
				\texttt{switch / case / otherwise / end} & Wybor jednej z wielu opcji \\
				\texttt{input("tekst")} & Wczytanie liczby od uzytkownika \\
				\texttt{input("tekst", "s")} & Wczytanie tekstu od uzytkownika \\
				\texttt{disp("komunikat")} & Wyswietlenie tekstu w oknie \\
				\texttt{plot(x, y, ...)} & Rysowanie wykresu \\
				\texttt{polyfit(x, y, n)} & Aproksymacja -- dopasowanie wielomianu \\
				\texttt{polyval(p, x)} & Obliczenie wartosci wielomianu \\
				\texttt{interp1(x, y, xi, m)} & Interpolacja -- wyznaczenie wartosci miedzy wezlami \\
				\bottomrule
			\end{tabular}
		\end{center}
	\end{infobox}
	
	\newpage
	% ============================================================
	\section*{\textcolor{headerblue}{BLOK 1 -- Decyzje: \texttt{if/else} i \texttt{switch/case}}}
	\noindent\textit{Rozpoczynamy od prostego kodu, ktory bedziemy stopniowo rozbudowywac.}
	% ============================================================
	
	% --- Cwiczenie 1 ---
	\begin{cwiczeniebox}
		\textbf{Cwiczenie 1 -- Szablon glowny} \quad \textit{(3 minuty)}
		
		Utworz szkielet programu, ktory bedzie rósł z kazdym cwiczeniem:
		
		\begin{enumerate}[leftmargin=*,itemsep=2pt]
			\item Wyczysc okno polecen (\texttt{clc}), zamknij wszystkie wykresy (\texttt{close all}) i usun zmienne z pamieci (\texttt{clear}).
			\item Zdefiniuj wezly $x$ i $y$ dla funkcji $f(x) = 2\sin(x) + x$:
			\begin{itemize}
				\item \texttt{x\_wezly = [0, 1, 2, 3, 4, 5, 6]}
				\item \texttt{y\_wezly = 2*sin(x\_wezly) + x\_wezly}
			\end{itemize}
		\end{enumerate}
	\end{cwiczeniebox}
	
	\newpage
	\begin{rozwiazaniebox}
		\textbf{Rozwiazanie Cwiczenia 1:}
		\begin{lstlisting}[caption={Cwiczenie 1 - Szablon glowny}]
			clear; close all; clc;
			
			x_wezly = [0, 1, 2, 3, 4, 5, 6];
			y_wezly = 2*sin(x_wezly) + x_wezly;
		\end{lstlisting}
	\end{rozwiazaniebox}
	
	\newpage
	% --- Cwiczenie 2 ---
	\begin{cwiczeniebox}
		\textbf{Cwiczenie 2 -- Dodajemy instrukcje warunkowe} \quad \textit{(5 minut)}
		
		Rozszerz program z Cwiczenia 1. Zapytajmy uzytkownika o badany punkt.
		
		\begin{enumerate}[leftmargin=*,itemsep=2pt]
			\item Zapytaj uzytkownika o dowolny punkt $x$ za pomoca \texttt{input} i przypisz go do zmiennej \texttt{x\_punkt}.
			\item Uzywajac instrukcji \texttt{if / elseif / else} sprawdz powiazanie tego punktu w stosunku do wezlow $x$ zdefiniowanych od 0 do 6:
			\begin{itemize}
				\item Jesli \texttt{x\_punkt} jest mniejsze od 0 – wyswietl \texttt{"Punkt poza wezlami - ekstrapolacja (lewa strona)."}
				\item Jesli \texttt{x\_punkt} jest wieksze od 6 – wyswietl \texttt{"Punkt poza wezlami - ekstrapolacja (prawa strona)."}
				\item W przeciwnym razie – wyswietl \texttt{"Punkt pomiedzy wezlami - interpolacja."}
			\end{itemize}
		\end{enumerate}
	\end{cwiczeniebox}
	
	\newpage
	\begin{rozwiazaniebox}
		\textbf{Rozwiazanie Cwiczenia 2:}
		\begin{lstlisting}[caption={Cwiczenie 2 - Instrukcje if/else}]
			x_punkt = input('Podaj badany punkt x: ');
			
			if x_punkt < 0
			disp('Punkt poza wezlami - ekstrapolacja (lewa strona).');
			elseif x_punkt > 6
			disp('Punkt poza wezlami - ekstrapolacja (prawa strona).');
			else
			disp('Punkt pomiedzy wezlami - interpolacja.');
			end
		\end{lstlisting}
	\end{rozwiazaniebox}
	
	\newpage
	% --- Cwiczenie 3 ---
	\begin{cwiczeniebox}
		\textbf{Cwiczenie 3 -- Dodajemy \texttt{switch/case}} \quad \textit{(5 minut)}
		
		Rozszerz program z Cwiczenia 2 przygotowujac wybor metody do pozniejszej interpolacji:
		
		\begin{enumerate}[leftmargin=*,itemsep=2pt]
			\item Za pomoca polecenia \texttt{disp} wyswietl \texttt{"Dostepne metody interpolacji: 1-linear, \\2-pchip, 3-spline"}.
			\item Zapytaj uzytkownika o metode (liczba 1, 2 lub 3) za pomoca \texttt{input}.
			\item Uzywajac instrukcji \texttt{switch / case / otherwise} stworz zmienna \texttt{metoda} (string) oraz \texttt{nazwa\_met} (string) zaleznie od wyboru:
			\begin{center}
				\begin{tabular}{ll}
					1 – \texttt{metoda='linear'}, \texttt{nazwa\_met='linear'} \\
					2 – \texttt{metoda='pchip'}, \texttt{nazwa\_met='pchip'} \\
					3 – \texttt{metoda='spline'}, \texttt{nazwa\_met='spline'}
				\end{tabular}
			\end{center}
			\item W przypadku zlej wartosci, wypisz \texttt{"Nieprawidlowy wybor -- domyslnie uzywam spline"} i przypisz powyzsze wartosci dla \texttt{spline}.
		\end{enumerate}
	\end{cwiczeniebox}
	
	\newpage
	\begin{rozwiazaniebox}
		\textbf{Rozwiazanie Cwiczenia 3:}
		\begin{lstlisting}[caption={Cwiczenie 3 - Instrukcja switch/case}]
			disp('Dostepne metody interpolacji: 1-linear, 2-pchip, 3-spline');
			wybor = input('Wybierz (1-3): ');
			
			switch wybor
			case 1
			metoda = 'linear'; nazwa_met = 'linear';
			case 2
			metoda = 'pchip'; nazwa_met = 'pchip';
			case 3
			metoda = 'spline'; nazwa_met = 'spline';
			otherwise
			disp('Nieprawidlowy wybor -- domyslnie uzywam spline');
			metoda = 'spline'; nazwa_met = 'spline';
			end
		\end{lstlisting}
	\end{rozwiazaniebox}
	
	\newpage
	% ============================================================
	\section*{\textcolor{headerblue}{BLOK 2 -- Wczytywanie danych i pierwsze wykresy}}
	% ============================================================
	
	% --- Cwiczenie 4 ---
	\begin{cwiczeniebox}
		\textbf{Cwiczenie 4 -- Dodajemy podstawowy wykres} \quad \textit{(8 minut)}
		
		Rozszerz program z Cwiczenia 3 o rysowanie wykresu:
		
		\begin{enumerate}[leftmargin=*,itemsep=2pt]
			\item Zapytaj uzytkownika: \texttt{"Czy narysowac wykres wezlow? (t/n): "} za pomoca \texttt{input}.
			\item Uzywajac \texttt{if} sprawdz odpowiedz.
			\item Wykres powinien spelniac nastepujace wymagania:
			\begin{itemize}
				\item Otwórz nowe okno \texttt{figure}
				\item Punkty wezlow narysuj jako czerwone kolka (\texttt{'ro'}) o rozmiarze 12 i grubosci 2.
				\item Dodaj siatke.
				\item Podpisz osie: \texttt{'x'} oraz \texttt{'y = 2*sin(x) + x')}.
				\item Dodaj tytul: \texttt{'Wezly funkcji f(x) = 2*sin(x) + x'}.
			\end{itemize}
		\end{enumerate}
	\end{cwiczeniebox}
	
	\newpage
	\begin{rozwiazaniebox}
		\textbf{Rozwiazanie Cwiczenia 4:}
		\begin{lstlisting}[caption={Cwiczenie 4 - Podstawowy wykres}]
			odp = input('Czy narysowac wykres wezlow? (t/n): ', 's');
			
			if odp == 't'
				figure('Name', 'Wykres wezlow', 'NumberTitle', 'off', ...
				'Position', [100, 100, 800, 600]);
				plot(x_wezly, y_wezly, 'ro', 'MarkerSize', 12, 'LineWidth', 2);
				grid on;
				xlabel('x', 'FontSize', 14);
				ylabel('y = 2*sin(x) + x', 'FontSize', 14);
				title('Wezly funkcji f(x) = 2*sin(x) + x', 'FontSize', 14);
			end
		\end{lstlisting}
	\end{rozwiazaniebox}
	
	\newpage
	% --- Cwiczenie 5 ---
	\begin{cwiczeniebox}
		\textbf{Cwiczenie 5 -- Dodajemy aproksymacje wielomianowa} \quad \textit{(10 minut)}
		
		Rozszerz program z Cwiczenia 4:
		
		\begin{enumerate}[leftmargin=*,itemsep=2pt]
			\item Zapytaj uzytkownika: \texttt{"Czy wykonac aproksymacje wielomianowa? (t/n): "}
			\item Po pozytywnej odpowiedzi:
			\begin{itemize}
				\item Zapytaj o stopien wielomianu (1-5).
				\item Zabezpiecz sie: jesli stopien < 1 ustaw na 2, jesli > 5 ustaw na 5.
				\item Uzyj \texttt{polyfit} do wyznaczenia wspolczynnikow.
				\item Utworz siatke \texttt{x\_siatka = 0 : 0.05 : 6} i uzyj \texttt{polyval}.
				\item Oblicz i wyswietl wynik aproksymacji dla zmiennej \texttt{x\_punkt} zapytanej w Cwiczeniu 2.
			\end{itemize}
			\item Narysuj wykres z wezlami, aproksymacja, funkcja oryginalna i nowo wyliczonym punktem.
		\end{enumerate}
	\end{cwiczeniebox}
	
	\newpage
	\begin{rozwiazaniebox}
		\textbf{Rozwiazanie Cwiczenia 5:}
		\begin{lstlisting}[caption={Cwiczenie 5 - Aproksymacja wielomianowa}]
			odp = input('Czy wykonac aproksymacje wielomianowa? (t/n): ', 's');
			
			if odp == 't'
			stopien = input('Podaj stopien wielomianu aproksymujacego (1-5): ');
			if stopien < 1
			disp('Za maly stopien -- ustawiono: 2'); stopien = 2;
			elseif stopien > 5
			disp('Za duzy stopien -- ustawiono: 5'); stopien = 5;
			else
			disp(['Wybrano stopien: ' num2str(stopien)]);
			end
			p = polyfit(x_wezly, y_wezly, stopien);
			disp('Wspolczynniki wielomianu (od najwyzszej potegi):');
			for i = 1:length(p), disp(['        p' num2str(i) ' = ' num2str(p(i))]); end
			
			x_siatka = 0 : 0.05 : 6;
			y_aproks = polyval(p, x_siatka);
			y_oryg = 2*sin(x_siatka) + x_siatka;
			
			y_punkt_aproks = polyval(p, x_punkt);
			disp(['Wartosc aproksymacji w punkcie x = ', num2str(x_punkt), ' wynosi y = ', num2str(y_punkt_aproks)]);
			
			figure('Name', 'Aproksymacja wielomianowa', 'NumberTitle', 'off', 'Position', [100, 100, 800, 600]);
			plot(x_wezly, y_wezly, 'ro', 'MarkerSize', 12, 'LineWidth', 2); hold on;
			plot(x_siatka, y_aproks, 'b-', 'LineWidth', 2);
			plot(x_siatka, y_oryg, 'g--', 'LineWidth', 1.5);
			plot(x_punkt, y_punkt_aproks, 'm*', 'MarkerSize', 15, 'LineWidth', 2);
			grid on;
			xlabel('x', 'FontSize', 14); ylabel('y', 'FontSize', 14);
			title(['Aproksymacja wielomianem stopnia ' num2str(stopien)], 'FontSize', 14);
			legend('Wezly pomiarowe', 'Wielomian aproksymujacy', 'Funkcja oryginalna', 'Wyliczony punkt', 'Location', 'best');
			end
		\end{lstlisting}
	\end{rozwiazaniebox}
	
	\newpage
	% ============================================================
	\section*{\textcolor{headerblue}{BLOK 3 -- Interpolacja}}
	% ============================================================
	
	% --- Cwiczenie 6 ---
	\begin{cwiczeniebox}
		\textbf{Cwiczenie 6 -- Dodajemy interpolacje} \quad \textit{(8 minut)}
		
		Rozszerz program z Cwiczenia 5 wykorzystujac wlasnosci zdefiniowane na samym poczatku:
		
		\begin{enumerate}[leftmargin=*,itemsep=2pt]
			\item Zapytaj uzytkownika:\\ \texttt{"Czy wykonac interpolacje wybrana wczesniej metoda? (t/n): "}
			\item Po pozytywnej odpowiedzi:
			\begin{itemize}
				\item Uzyj zmiennej \texttt{metoda} z Cwiczenia 3 w funkcji \texttt{interp1} do interpolacji na siatce \texttt{0:0.05:6}.
				\item Oblicz i wyswietl wynik interpolacji uzywajac wczesniej wpisanego \texttt{x\_punkt}\\(z Cwiczenia 2).
			\end{itemize}
			\item Narysuj wykres z wezlami, interpolacja, funkcja oryginalna i nowo wyliczonym punktem.
		\end{enumerate}
	\end{cwiczeniebox}
	
	\newpage
	\begin{rozwiazaniebox}
		\textbf{Rozwiazanie Cwiczenia 6:}
		\begin{lstlisting}[caption={Cwiczenie 6 - Interpolacja}]
			odp = input('Czy wykonac interpolacje wybrana wczesniej metoda? (t/n): ', 's');
			
			if odp == 't'
			x_siatka = 0 : 0.05 : 6;
			y_interp = interp1(x_wezly, y_wezly, x_siatka, metoda);
			y_oryg = 2*sin(x_siatka) + x_siatka;
			
			y_punkt_interp = interp1(x_wezly, y_wezly, x_punkt, metoda);
			disp(['Wartosc interpolacji w punkcie x = ', num2str(x_punkt), ' wynosi y = ', num2str(y_punkt_interp)]);
			
			figure('Name', 'Interpolacja', 'NumberTitle', 'off', 'Position', [100, 100, 800, 600]);
			plot(x_wezly, y_wezly, 'ro', 'MarkerSize', 12, 'LineWidth', 2); hold on;
			plot(x_siatka, y_interp, 'b-', 'LineWidth', 2);
			plot(x_siatka, y_oryg, 'g--', 'LineWidth', 1.5);
			plot(x_punkt, y_punkt_interp, 'm*', 'MarkerSize', 15, 'LineWidth', 2);
			
			grid on;
			xlabel('x', 'FontSize', 14); ylabel('y', 'FontSize', 14);
			title(['Interpolacja metoda ' nazwa_met], 'FontSize', 14);
			legend('Wezly pomiarowe', 'Krzywa interpolacji', 'Funkcja oryginalna', 'Wyliczony punkt', 'Location', 'best');
			end
		\end{lstlisting}
	\end{rozwiazaniebox}
	
	\newpage
	% ============================================================
	\section*{\textcolor{headerblue}{Podsumowanie}}
	% ============================================================
	
	\bigskip
	\begin{infobox}
		\textbf{Kluczowe roznice do zapamietania:}
		\begin{itemize}
			\item \texttt{if} sprawdza \textbf{dowolny warunek logiczny}; \texttt{switch} porownuje wartosc do \textbf{konkretnych opcji}.
			\item \texttt{input("pytanie")} zwraca \textbf{liczbe}; \texttt{input("pytanie", "s")} zwraca \textbf{tekst}.
			\item \textbf{Aproksymacja}: krzywa \textbf{nie musi} przechodzic przez wezly – szukamy krzywej najlepiej dopasowanej.
			\item \textbf{Interpolacja}: krzywa \textbf{zawsze} przechodzi przez kazdy wezel.
			\item \textbf{Metody interpolacji}: \texttt{linear} (lamana), \texttt{pchip} (bez oscylacji), \texttt{spline} (najgladsza).
			\item \texttt{pause} wstrzymuje program – uzywamy go miedzy etapami dla lepszej czytelnosci.
		\end{itemize}
	\end{infobox}
	
\end{document}